%\ifodd\value{page}\hbox{}\newpage\fi
\chapter*{Śrī Lalitā Sahasranāma — Śloka 14}
\addcontentsline{toc}{chapter}{Śloka 14 - Nomi 30,31}

\begin{sloka}
{\sanskritfont
\begin{tabular}{c}
कामेश्वरप्रेमरत्नमणिप्रतिपणस्तनी ।\\
नाभ्यालवालरोमालिलताफलकुचद्वयी ॥ १४॥
\end{tabular}

\vspace{1.5em}

{\middlesize \iastfont
kāmeśvara-prema-ratna-maṇi-prati-paṇa-stanī |\\
nābhy-ālavāla-romāli-latā-phala-kuca-dvayī ||14||}

\vspace{1em}

{\middlesize
\anvaya{%
कामेश्वरप्रेमरत्नमणिप्रतिपणस्तनी ।
नाभ्यालवालरोमालिलताफलकुचद्वयी ।
}
}

\middlesize \iastfont \itmean{%
«I cui seni sono offerti come pegno prezioso, simili a gemme, nell’amore di Kāmeśvara;  
i cui due seni sono come frutti portati da una liana che sorge dal bacino dell’ombelico.»%
}

}
\end{sloka}

\wordmeaning{%
\wordline{कामेश्वर-प्रेम-रत्न-मणि-प्रति-पण-स्तनी}
{kāmeśvara-prema-ratna-maṇi-prati-\\paṇa-stanī}
{ i cui seni (\textit{stanī}) sono un pegno o offerta (\textit{prati-paṇa}), paragonati a gemme preziose (\textit{ratna-maṇi}), all’interno dell’amore (\textit{prema}) di \textit{Kāmeśvara}; }%

\wordline{नाभि-आलवाल-रोम-आलि-लता-फल-कुच-द्वयी}
{nābhi-ālavāla-roma-āli-latā-\\phala-kuca-dvayī}
{ i cui due seni (\textit{kuca-dvayī}) sono paragonati a frutti (\textit{phala}) portati da una liana (\textit{latā}) di sottili peli (\textit{roma-āli}) che sorge dal bacino dell’ombelico (\textit{nābhi-ālavāla}); }%
}

\textit{Śloka} 14 prosegue il movimento discendente della contemplazione, soffermandosi sui seni e sull’ombelico come un unico campo simbolico. I seni sono anzitutto descritti come un pegno prezioso, liberamente offerto nell’amore reciproco tra \textit{Lalitā} e \textit{Kāmeśvara}. Il desiderio qui non è possesso, ma dono consacrato.

Nel secondo \textit{pāda}, l’ombelico è immaginato come un bacino fertile da cui sorge una liana sottile, che porta frutto nei seni stessi. L’immagine è organica e continua: nutrimento, crescita e bellezza appaiono come aspetti di una medesima forma vivente.

Nel linguaggio simbolico dello \textit{Śrī Lalitā Sahasranāma}, i seni non sono descritti come elementi erotici né come oggetti di sguardo, ma come segni visibili della potenza che nutre e sostiene la vita. Essi rappresentano la capacità generativa e dispensatrice della \textit{Śakti}: ciò da cui il mondo riceve nutrimento, continuità e protezione. La loro pienezza non allude alla seduzione, ma all’abbondanza; non alla possessività, ma al dono. Quando sono detti “pegno nell’amore di \textit{Kāmeśvara}”, indicano un’offerta reciproca e consacrata tra coscienza e potenza vitale, non un atto di appropriazione. In questa prospettiva, il corpo della Dea diventa una mappa cosmica: ciò che appare come forma è in realtà funzione, e ciò che è descritto come bellezza è espressione di un’energia che sostiene l’esistenza stessa.


Questo \textit{śloka} contiene due \textit{nāma}.  
\textbf{\textit{Kāmeśvara-prema-ratna-maṇi-prati-paṇa-stanī}} presenta i seni come offerte preziose nell’amore divino.  
\textbf{\textit{Nābhy-ālavāla-romāli-latā-phala-kuca-dvayī}} raffigura l’ombelico e i seni come una continuità fertile e vivente della forma.
